\Needspace{18\baselineskip}
\section{Original-hole tuples and concrete safety predicates}
\label{sec:tuple-safety}
This section discharges the former third item in
Appendix~\ref{app:remaining-work}: compatibility of recursive specimens with
finite tuple substitution, and definitions of the concrete quasi-specimen,
input-safety, and robustness predicates in paper \S\S8.1--8.3.
The compatibility theorem concerns the existing finite-depth representation.
The separate comparison with set-based hypernets and changes of depth bound
is now proved in Section~\ref{sec:paper-encoding}.

The tuple semantics uses simultaneous graph expansion at each original graph
layer. Its recursive calls visit original box bodies; a hole replacement is a
leaf, never a fresh context to be traversed. This is a direct semantics for
parallel application at named holes, with retained edges represented by
one-edge copies. It is not a new claim about arbitrary sequential schedules.
The earlier ordinary-substitution and interchange theorems concern those
schedules. Here the substantive checks are that the tuple is finite and
aligned, that its names are distinct original holes, that no original hole is
omitted when the result is closed, and that both representations produce the
same filled graphs.

The concrete safety predicates below are definitions, not proofs that arbitrary
pre-templates satisfy them. They retain the clauses that the abstract interface
specialized away. Two lemmas then derive the rooted consequences needed by the
abstract proof. Constructing that interface, its full-UAM context-step law,
and the final sufficiency instantiation remains unfinished.

\Needspace{18\baselineskip}
\subsection{Finite tuples and their application}
\termdef{tuple-entry}
\begin{definition}[Aligned substitution entry]
An entry records a depth $r$, ordered interfaces $I,O$, a hole name $h$,
equalities $I=h.\mathsf{input}$ and $O=h.\mathsf{output}$, and two graphs
$H,K:I\to O$ at depth $r$. A finite list of entries therefore represents the
three equally long vectors of names, left replacements, and right replacements.
The varying depth and interface of an entry are stored with its two graphs.
\end{definition}
The constructor at an actual primitive hole reads its interface from the
selected edge. Rewriting its label to the hole label proves the two interface
equalities. An entry is related when its two graphs satisfy the pre-template
at its recorded depth and interface; this is exactly the pointwise condition
on the two vectors.
\lean{Graph/TupleFillings}{9}{17}{Linear.TupleEntry}
\lean{Graph/TupleFillings}{20}{31}{Linear.TupleEntry.at}
\lean{Graph/TupleFillings}{34}{35}{Linear.TupleEntry.Related}

\begin{definition}[Distinct traversal enumeration]
Remove repetitions from a finite enumeration by recursively processing the
tail and retaining the head only when it is absent from the processed tail.
Apply this operation to the graph's finite edge enumeration.
\end{definition}
\begin{lemma}[Enumeration laws]
Removing repetitions preserves membership and gives a list with no duplicates.
Consequently every original shallow edge occurs exactly once in the traversal.
\end{lemma}
\begin{proof}
Induct on the input list. If the head is already in the processed tail, removing
it does not change membership; otherwise prepend it. The same cases preserve
no-duplicates, using the tested absence in the second case. Completeness of
the graph's supplied enumeration now gives membership of each edge in the
distinct traversal list. This matters because the underlying finite-carrier
structure permits repeated enumeration entries.
\end{proof}
\lean{Graph/TupleFillings}{38}{42}{Linear.tupleDistinct}
\lean{Graph/TupleFillings}{44}{56}{Linear.mem_tupleDistinct}
\lean{Graph/TupleFillings}{58}{66}{Linear.tupleDistinct_nodup}
\lean{Graph/TupleFillings}{69}{70}{Linear.tupleEdges}
\lean{Graph/TupleFillings}{72}{75}{Linear.mem_tupleEdges}

\termdef{tuple-filling}
\begin{definition}[Tuple application in traversal order]
A tuple-filling certificate simultaneously describes two filled graphs and an
aligned list of entries. At a graph $C$, choose a left and right replacement
for each shallow edge, expand both families, and concatenate the local entry
lists in the distinct edge traversal order. Each original edge has one of
three local certificates:
\begin{enumerate}
\item Retain its label and ordered boundary as one-edge copies on both sides;
its entry list is empty.
\item At an actual hole, insert a pointwise related pair and record the one
corresponding entry. Do not recurse into either inserted graph.
\item At a box, recursively apply tuples to the original body, retain its
binder metadata and external interface, and install one-edge copies of the
boxes with those two resulting bodies. Its entry list is the body's list.
\end{enumerate}
\end{definition}
The clauses allow partial filling of an open context. The completeness theorem
below proves that a closed result forces all original holes to have been
filled. Thus the general certificate is also useful before the state
restriction is imposed.
\lean{Graph/TupleFillings}{79}{86}{Linear.TupleFilling}
\lean{Graph/TupleFillings}{89}{109}{Linear.TupleFillingEdge}

\begin{theorem}[Compatibility with recursive filling]
There exists a finite tuple-filling certificate from $C$ to $L,R$ if and only
if there is a recursive deep-filling certificate with those same three graphs.
\end{theorem}
\begin{proof}
For the forward direction, forget the entry lists and induct on the depth.
The retained and filled clauses give their original certificates directly;
apply the induction hypothesis to a boxed body. For the reverse direction,
induct on depth again. At each edge record the empty list, the singleton
replacement entry, or the recursively obtained body's list. Choose those
local witnesses and concatenate them over the finite distinct traversal.
The two expansions have exactly the original carriers, incidences, and labels,
so neither endpoint graph is changed by extracting the tuple.
\end{proof}
\lean{Graph/TupleFillings}{112}{136}{Linear.TupleFilling.forget}
\lean{Graph/TupleFillings}{139}{162}{Linear.DeepFilling.tuple}
\lean{Graph/TupleFillings}{165}{168}{Linear.tupleFilling_iff}

\begin{lemma}[Pointwise relatedness]
Every recorded entry satisfies the pre-template relation at its own depth.
\end{lemma}
\begin{proof}
Membership in a concatenated list selects one original edge and one entry in
its local list. A retained edge contributes no entries. A filled edge carries
the required relation. For a box use induction on its original body's depth.
\end{proof}
\lean{Graph/TupleFillings}{171}{201}{Linear.TupleFilling.related}

\subsection{Exact original-hole inventory}
\begin{lemma}[Provenance]
Every recorded name is a hole of the original context. More precisely, an
entry in the local list for edge $e$ comes from a hole in the occurrence fibre
owned by $e$.
\end{lemma}
\begin{proof}
A filled edge is itself the named hole. A box's entries come, by induction,
from holes in its original body; these are exactly its child occurrences in
the outer graph. A retained edge contributes nothing. The graph-level result
selects the owning edge using membership in the concatenation. No step of
this argument looks for a name inside a replacement graph.
\end{proof}
\lean{Graph/TupleFillings}{204}{239}{Linear.TupleFilling.provenance}
\lean{Graph/TupleFillings}{242}{255}{Linear.TupleFillingEdge.provenance}

\begin{theorem}[Distinct tuple names]
If the original context has unique hole labels, the list of tuple names has
no duplicates.
\end{theorem}
\begin{proof}
Within one local list, the retained and filled cases are empty and singleton.
For a box, uniqueness of holes in the original outer context implies
uniqueness in its body; apply induction there. Between different original
edges, two entries with the same name would give occurrences with that name
in both fibres, by local provenance. Original hole uniqueness identifies the
occurrences and hence their owning edges, a contradiction. Finally the edge
traversal itself has no duplicates. These within-list and between-list facts
prove pairwise distinctness of the concatenation and therefore of its names.
\end{proof}
\lean{Graph/TupleFillings}{258}{298}{Linear.TupleFilling.unique}

\begin{lemma}[Reflecting closedness]
If a simultaneous expansion is hole-free, every graph inserted by that
expansion is hole-free. A hole-free one-edge copy of a box has a hole-free body.
\end{lemma}
\begin{proof}
An occurrence of a hole in an inserted graph lifts to the expansion by tagging
its outer edge with the owning original edge; the label-local occurrence is
unchanged. This contradicts closedness. For a copied box, a body hole lifts to
a child occurrence of its unique edge, using the equality of the box label.
\end{proof}
\lean{Graph/TupleFillings}{301}{307}{Linear.tuple_component_closed}
\lean{Graph/TupleFillings}{310}{315}{Linear.LabelCopy.body_closed}

\begin{theorem}[Completeness of the tuple]
If the left filled graph is hole-free, every original hole name occurs in
the tuple. If the context also has unique holes, its holes are enumerated
exactly once by the tuple names.
\end{theorem}
\begin{proof}
Choose an occurrence of the original hole and its owning edge. If that edge
were retained, its one-edge copy would retain the same hole occurrence,
contradicting the reflected closedness of its inserted graph. If the edge was
filled, the primitive hole label has only its shallow occurrence, so the
recorded name is exactly the given name. If it is a box, reflect closedness to
the resulting body and apply induction to the original body occurrence.
Membership of the owning edge in the traversal puts the entry in the full
list. Combining completeness, provenance, and distinctness gives the exact
inventory statement.
\end{proof}
\lean{Graph/TupleFillings}{318}{376}{Linear.TupleFilling.complete}
\lean{Graph/TupleFillings}{380}{394}{Linear.TupleFilling.inventory}
For an actual recursive specimen, its context already has unique holes and
both expanded graphs are closed. Tuple extraction therefore supplies the same
left and right state graphs, pointwise related entries, distinct names, and
the exact original-hole inventory.
\lean{Graph/TupleFillings}{398}{405}{Linear.DeepSpecimen.tuple}

\begin{lemma}[Closed replacement entries]
When both filled results are hole-free, both graphs in every recorded entry
are hole-free.
\end{lemma}
\begin{proof}
Reflect closedness through expansion to each selected local graph. In the
filled case these are exactly the entry's two replacements. In the box case
reflect through the two copied box labels to their bodies, then apply
induction. There are no entries in a retained case.
\end{proof}
This also rules out holes hidden inside replacement graphs in a successful
state-producing tuple; they cannot have disappeared through later recursive
visits, because such visits are absent from the application clauses.
\lean{Graph/TupleFillings}{408}{441}{Linear.TupleFilling.entries_closed}

\termdef{named-tuple-application}
\begin{definition}[Named tuple application]
An arbitrarily ordered aligned tuple applies when it is a permutation of a
traversal-order tuple-filling certificate with the required endpoints.
All components of an entry move together. Naming by distinct original holes
makes the order of the displayed vectors immaterial.
\end{definition}
Existence of such an application is equivalent to recursive filling: erase
the permutation witness and forget the tuple in one direction; use the
identity permutation on an extracted tuple in the other. Permuting an
application again preserves it by composition of list permutations. This is
order invariance of the parallel named-tuple presentation, not a separate
sequential-scheduling theorem.
\lean{Graph/TupleFillings}{445}{447}{Linear.AppliesTuple}
\lean{Graph/TupleFillings}{450}{456}{Linear.appliesTuple_iff}
\lean{Graph/TupleFillings}{459}{463}{Linear.AppliesTuple.permute}

\termdef{specimen-tuple}
\begin{definition}[Complete tuple of a specimen]
A specimen tuple packages its entry list, named application to its actual
filled state graphs, pointwise relatedness, closed replacements, distinct
names, and exact coverage of the context's original holes.
\end{definition}
Every recursive specimen supplies this package: choose its extracted list,
use identity permutation for application, and apply the relatedness,
closedness, and inventory theorems. No compatibility assertion is postulated
as a field without constructing its proof.
\lean{Graph/TupleFillings}{466}{473}{Linear.SpecimenTuple}
\lean{Graph/TupleFillings}{476}{481}{Linear.DeepSpecimen.substitutionTuple}

\begin{theorem}[Reconstructing a specimen]
Conversely, tuple application in an allowed focussed unique-hole context,
with closed left and right results, constructs a recursive specimen with that
context, the same focus mode, and exactly those two result graphs.
\end{theorem}
\begin{proof}
Forget each local tuple certificate, using the already proved graph-level
forgetting theorem for boxed bodies. The outer tuple constructor supplies its
two replacement families. Put those families and the forgotten local
certificates in the specimen record, together with the given context, focus,
allowedness, and closedness proofs. Its left and right state graph projections
are definitionally the supplied expansions. For an arbitrarily ordered named
tuple, first take its traversal-order application witness.
\end{proof}
\lean{Graph/TupleFillings}{484}{490}{Linear.TupleFillingEdge.forget}
\lean{Graph/TupleFillings}{494}{513}{Linear.TupleFilling.specimen}

\Needspace{18\baselineskip}
\subsection{The full quasi-specimen predicate}
In the remaining definitions, fix a layered pre-template $T$, an allowed
class $\mathbb C$ of erased focus-free contexts, an operation partition, and
an admissible behaviour $B$. Let $M$ be the concrete full UAM of
Section~\ref{sec:full-uam}. All executions below are executions of $M$, not
search-only executions. A concrete specimen is the existing recursive
specimen, whose tuple interpretation was just established.

\begin{definition}[Settled and remote concrete specimens]
A specimen is settled when its context focus is backtracking, or searching
without entering a hole. It is remote when its focus is not a rewrite focus
or is not entering a hole. Settled implies remote, by distinguishing the two
settled modes from rewrite mode.
\end{definition}
\lean{Graph/ConcreteSafety}{14}{15}{Linear.ConcreteSafety.Settled}
\lean{Graph/ConcreteSafety}{18}{19}{Linear.ConcreteSafety.Remote}
\lean{Graph/ConcreteSafety}{113}{116}{Linear.ConcreteSafety.settled_remote}

\termdef{concrete-quasi}
\begin{definition}[Concrete quasi-specimen]
For relations $A,D$ on actual states, a pair $s,t$ is a quasi-specimen up to
$(A,D)$ when there exists an allowed specimen $S$ satisfying:
\begin{enumerate}
\item The focus modes of $S$, $s$, and $t$ agree.
\item If $s$ and $t$ are both rooted, then both filled states $S_L,S_R$ are
rooted, $A(s,S_L)$, and $D(S_R,t)$.
\end{enumerate}
The existence of $S$ and the mode equalities are unconditional. Only the
second clause is conditional on rootedness. A non-rewrite quasi-specimen
additionally requires that the common mode is not rewrite.
\end{definition}
These are the clauses of paper \S8.1. In particular, an unrooted pair does
not satisfy the definition merely by making a rootedness implication vacuous:
a mode-compatible specimen must still exist.
\lean{Graph/ConcreteSafety}{23}{28}{Linear.ConcreteSafety.Quasi}
\lean{Graph/ConcreteSafety}{31}{32}{Linear.ConcreteSafety.NonRewriteQuasi}

\begin{lemma}[Elementary quasi-specimen laws]
Every actual specimen yields a quasi-specimen of its two filled states up to
equality. Every quasi-specimen has equal endpoint modes. Weakening either
up-to relation preserves quasi-specimens.
\end{lemma}
\begin{proof}
For the first claim use the specimen itself: focus preservation gives the
mode equalities, and the conditional conclusions are the supplied rootedness
hypotheses and reflexive equalities. For the second, compose the two mode
equalities through the witness's mode. For weakening, retain that same
witness and compose its two conditional comparison facts with the supplied
implications of relations.
\end{proof}
\lean{Graph/ConcreteSafety}{35}{37}{Linear.ConcreteSafety.specimen_quasi}
\lean{Graph/ConcreteSafety}{40}{43}{Linear.ConcreteSafety.Quasi.modes}
\lean{Graph/ConcreteSafety}{46}{52}{Linear.ConcreteSafety.Quasi.mono}

\begin{lemma}[Rooted non-rewrite consequence]
Given a concrete non-rewrite quasi-specimen and rooted endpoints, there is a
specimen whose mode is not rewrite, whose two filled states are rooted, and
which provides the two up-to comparisons.
\end{lemma}
\begin{proof}
Take the unconditional witness. Its mode equality transfers the endpoint's
non-rewrite property to it. Apply its conditional clause to the two rootedness
hypotheses. This derives the consequence used by the abstract proof; it does
not identify the stronger concrete definition with the weaker abstract one.
\end{proof}
\lean{Graph/ConcreteSafety}{56}{65}{Linear.ConcreteSafety.NonRewriteQuasi.rooted_witness}

\Needspace{18\baselineskip}
\subsection{Input-safety and robustness with the paper's clauses}
\termdef{concrete-input-safety}
\begin{definition}[Concrete input-safety]
For every allowed specimen $S$ with an entering search focus, at least one of
the following holds. There is no rootedness restriction on $S$.
\begin{enumerate}
\item Both filled states reach stuck states in finitely many full-UAM steps.
\item For some settled specimen $D$ and $k,l\in\mathbb N$,
$(1+k)\,Q\,l$, $S_L\steps{1+k}D_L$, and $S_R\steps{l}D_R$.
\item For some $k,l\in\mathbb N$ and non-rewrite quasi-specimen $s,t$ up to
$(\simeq_P,\simeq_P)$,
$(1+k)\,Q\,(1+l)$, $S_L\steps{1+k}s$, and $S_R\steps{1+l}t$.
\end{enumerate}
\end{definition}
Here stuckness and the state equivalences use $M$. In the third alternative
both up-to relations are equivalences, as in the paper, not only refinements.
\lean{Graph/ConcreteSafety}{68}{77}{Linear.ConcreteSafety.InputSafe}

\termdef{concrete-robustness}
\begin{definition}[Concrete robustness relative to a rewrite]
For a rewrite pair $x\to x'$, consider every allowed specimen $S$ with
$S_L=x$ and a non-entering rewrite focus. At least one of these holds:
\begin{enumerate}
\item $S_L$ or $S_R$ is not rooted.
\item $x'$ reaches a stuck state in finitely many full-UAM steps.
\item For some $k,l\in\mathbb N$ and non-rewrite quasi-specimen $s,t$ up to
$(\preceq_P,\preceq_R)$,
$(1+k)\,Q\,l$, $x'\steps{k}s$, and $S_R\steps{l}t$.
\end{enumerate}
Robustness for the UAM requires this property for every substitution or
extrinsic rewrite transition of the supplied behaviour.
\end{definition}
The relative predicate is stated on a pair; the global predicate restricts
its quantified pairs to actual rewrites. This preserves the paper's unrooted
alternative and its asymmetric step counts, including the one already taken
on the left.
\lean{Graph/ConcreteSafety}{81}{88}{Linear.ConcreteSafety.RobustRelative}
\lean{Graph/ConcreteSafety}{91}{92}{Linear.ConcreteSafety.Robust}
\termdef{concrete-template}
\begin{definition}[Concrete template]
A template is input-safe and output-closed. In the depth-indexed presentation,
output closure is required at every available layer of the pre-template.
\end{definition}
\lean{Graph/ConcreteSafety}{95}{96}{Linear.ConcreteSafety.Template}

\begin{lemma}[Restricting robustness to rooted specimens]
For a robust pre-template and a non-entering rewrite specimen with two rooted
filled states, every rewrite of its left state has either the stuck-tail
alternative or the matching non-rewrite quasi-specimen alternative above.
\end{lemma}
\begin{proof}
Apply concrete robustness to that actual rewrite and specimen. Its first
alternative contradicts one of the two supplied rootedness hypotheses; the
remaining disjunction is exactly the claimed result. This is the legitimate
reason the abstract rooted interface can omit the unrooted alternative.
\end{proof}
\lean{Graph/ConcreteSafety}{100}{110}{Linear.ConcreteSafety.Robust.rooted}

\subsection{Regression examples and scope}
The existing example has one hole strictly inside a box and no shallow hole.
Its complete substitution tuple contains that deep hole. In fact its list of
names is exactly the singleton containing that name: completeness proves
membership, provenance and the example's occurrence classification exclude
any other name, and no-duplicates excludes a second entry. Every recorded
replacement is closed by the general reflected-closedness theorem. These
checks exercise genuine recursion through a box, not just shallow filling.
\lean{Graph/TupleSafetyExamples}{9}{10}{Linear.TupleSafetyExamples.nested_hole_recorded}
\lean{Graph/TupleSafetyExamples}{13}{37}{Linear.TupleSafetyExamples.nested_tuple_singleton}
\lean{Graph/TupleSafetyExamples}{40}{41}{Linear.TupleSafetyExamples.nested_entries_closed}
The same example gives a quasi-specimen up to equality although at least one
of its filled states is not rooted, by the previously proved path obstruction.
Thus the definition has not accidentally required unconditional rootedness.
Conversely, an initial search-mode value and a final backtrack-mode value
cannot form a quasi-specimen, even if both up-to relations are universally
true: the unconditional endpoint-mode equality would equate two distinct
focus constructors.
\lean{Graph/TupleSafetyExamples}{44}{47}{Linear.TupleSafetyExamples.unrooted_quasi}
\lean{Graph/TupleSafetyExamples}{51}{59}{Linear.TupleSafetyExamples.different_modes_not_quasi}

The tuple and concrete-definition obligation is removed from
Appendix~\ref{app:remaining-work}. This does not assert that a particular
nontrivial pre-template is input-safe or robust. Section~\ref{sec:concrete-specimen-system}
now constructs the named-state specimen interface and proves the safety and
robustness bridges. General initial observation lifting and state-class
integration still separate these results from the final theorem.
